\section{Introduction and research question}
Hydrogen is an industrial feedstock used in applications such as refining and chemical manufacture. A common production route is steam methane reforming (SMR): natural gas reacts with steam to form a hydrogen-rich synthesis gas. The chemistry and the high-temperature reformer both create a carbon challenge. Carbon entering with methane ultimately appears largely as CO$_2$, while conventional reformers also burn fuel to supply reaction heat.

Carbon capture and storage (CCS) can separate and store much of the concentrated process CO$_2$, but it does not by itself eliminate fired reformer heat or upstream natural-gas emissions. This motivates the question tested here: can high-temperature nuclear heat replace the fired heat service while CCS manages process carbon?

A high-temperature gas-cooled reactor (HTGR) is relevant because helium coolant can deliver temperatures suitable for industrial process heat. An intermediate heat exchanger (IHX) isolates the nuclear primary-helium circuit from a separate secondary-helium circuit feeding the reformer. The proposal therefore transfers heat across an engineered boundary rather than sending reactor primary coolant into the chemical process.

For Singapore, this remains a future-oriented screening question rather than an existing deployment option. Nuclear licensing/siting and cross-border CO$_2$ storage are not established project infrastructure. CN4252 nevertheless supplies a quantitative test: the proposed solution must avoid more than 0.25 MtCO$_2$e/y at an abatement cost below S\$100/tCO$_2$e.

The research question is therefore: \emph{does the strongest directly supported HTGR-assisted SMR+CCS design found in the primary literature satisfy both CN4252 thresholds when its process balance is extended to a transparent Singapore lifecycle and forward-economic boundary?}

The final process basis is INL's NGNP HTGR-integrated SMR analysis~\cite{inltev953,inltev961}; Nishihara/JAEA evidence supplies the selected HTGR/IHX hardware and mature design-study economic architecture~\cite{nishihara2007potential}. Historical project configurations remain in the repository for audit but are not the submission-facing design.
